\documentclass[12pt,a4paper]{amsart}

\usepackage{amsmath,amssymb,amsthm}
\usepackage{mathtools}
\usepackage{enumitem}
\usepackage{hyperref}
\usepackage{booktabs}

%%%─── Theorem environments
\newtheorem{theorem}{Theorem}[section]
\newtheorem{lemma}[theorem]{Lemma}
\newtheorem{proposition}[theorem]{Proposition}
\newtheorem{corollary}[theorem]{Corollary}
\newtheorem{conjecture}[theorem]{Conjecture}
\newtheorem{hypothesis}[theorem]{Hypothesis}
\theoremstyle{definition}
\newtheorem{definition}[theorem]{Definition}
\newtheorem{assumption}[theorem]{Assumption}
\newtheorem{problem}[theorem]{Open Problem}
\theoremstyle{remark}
\newtheorem{remark}[theorem]{Remark}
\newtheorem*{remark*}{Remark}

%%%─── Macros
\newcommand{\norm}[1]{\|#1\|}
\newcommand{\ip}[2]{\langle #1,\, #2\rangle}
\newcommand{\Kc}{\mathcal{K}_c}
\newcommand{\lam}[2]{\lambda_{#2}^{(#1)}}
\newcommand{\gap}[2]{\mathrm{gap}_{#2}^{(#1)}}
\newcommand{\pswf}[1]{\psi_{#1}^{(c)}}
\newcommand{\op}{\mathrm{op}}
\newcommand{\spann}{\mathrm{span}}

\subjclass[2020]{47B34, 45C05, 35P20, 47A55, 81Q10}
\keywords{Prolate spheroidal wave functions, sinc kernel, spectral gaps,
  Gram coercivity, Airy scaling, quasimodes, compact self-adjoint operators,
  Hilbert--P\'{o}lya programme}

\begin{document}

\title[Gap-S: Reduction Framework]{%
  Paper~XIII: A Reduction Framework for Gap-S\\[0.3em]
  \large Gram Coercivity as Stability Input,
  Quasimode Transfer as Gap Mechanism}

\author{\textit{prolate-gram-coercivity programme}}
\date{May 2026}

\begin{abstract}
This paper establishes a reduction framework for the spectral gap problem
(Gap-S) for the prolate concentration operator
\[
  (\mathcal{K}_c f)(x) := \int_{-1}^1
  \frac{\sin(c(x-y))}{\pi(x-y)}\,f(y)\,dy
\]
in the transition regime $n \sim 2c/\pi$.

\medskip\noindent
\textbf{Central negative result.}
Gram coercivity ($\lambda_{\min}(G_N) \asymp c^{-1/2}$, Paper~I)
does \emph{not} imply nontrivial spectral gaps
(Proposition~\ref{prop:gram-fails}).
The example $T_E = \lambda I$ shows this concretely.
Gram coercivity is a \emph{stability input}, preventing coordinate collapse,
but carries no eigenvalue repulsion.

\medskip\noindent
\textbf{Central structural result.}
For $\Kc$-invariant trial spaces, all ingredients needed to force
Gap-S reduce circularly to Gap-S itself
(Remark~\ref{rem:circular}).
This explains why Papers~X--XII could not close the argument,
and identifies Option~B of Paper~XII
(non-$\Kc$-invariant phase-space trial spaces) as the only non-circular route.

\medskip\noindent
\textbf{Main outputs.}
\begin{itemize}
\item \emph{(Theorem~A, Compressed Spectral Transfer.)}
  Under Gram coercivity, local spectral nondegeneracy
  (Assumption~\ref{ass:nondeg}), and spectral window isolation
  (Assumption~\ref{ass:window}),
  one obtains $\gap{c}{n} \geq Cc^{-1/2}$.
  Theorem~A is a \emph{lifting mechanism}:
  Assumption~\ref{ass:nondeg} encodes most of the gap information;
  the theorem's value lies in isolating what is needed.
\item \emph{(Theorem~B, Quasimode Transfer.)}
  Under local quasimode hypotheses (Assumption~\ref{ass:quasimode}),
  a sharp conditional bound at model scale follows.
  This is the most solid result in the paper.
  Condition~(Q5) is the genuinely hard spectral counting assumption;
  it encodes most of the transition-window spectral rigidity.
\item \emph{(Corollary~C.)} Under Airy transition asymptotics,
  the gap scale would be $c^{-1/3}$.
  This is a \emph{target theorem}, conditional on a
  fully justified Airy quasimode package not yet constructed.
\end{itemize}
\end{abstract}

\maketitle
\tableofcontents
\bigskip\hrule\bigskip

%% ============================================================
\section{Context and Position in the Programme}
\label{sec:context}
%% ============================================================

\begin{itemize}[nosep]
  \item \textbf{Paper~I}: Gram coercivity $\lambda_{\min}(G_N^{(c)}) \asymp c^{-1/2}$
    --- proved unconditionally.
  \item \textbf{Paper~IX}: B-strong $\Leftrightarrow$ $\varepsilon_N^{(c)} \geq Cc^{-1/2}$.
  \item \textbf{Papers~X--XII}: All direct routes fail; the obstruction is structural.
  \item \textbf{This paper}: Gram coercivity is a stability input only.
    The minimal ingredients for Gap-S are isolated;
    the circularity of invariant-space approaches is proved.
\end{itemize}

\subsection{Gap-S}

Gap-S asks for $g(c) > 0$ with
\begin{equation}\label{eq:gapS}
  |\lam{c}{n} - \lam{c}{n+1}| \geq g(c) \gg c^{-1/4},
  \quad n \sim 2c/\pi.
\end{equation}

%% ============================================================
\section{Operator Setup}
\label{sec:setup}
%% ============================================================

\begin{definition}[Prolate concentration operator]
$(\Kc f)(x) := \int_{-1}^1 \tfrac{\sin(c(x-y))}{\pi(x-y)} f(y)\,dy$;
eigenvalues $1 > \lam{c}{0} \geq \lam{c}{1} \geq \cdots > 0$.
\end{definition}

\begin{definition}
$\gap{c}{n} := \lam{c}{n} - \lam{c}{n+1}$,\quad
$V_N := \spann\{\pswf{0},\ldots,\pswf{N}\}$.
\end{definition}

\begin{lemma}[Compression --- Rayleigh quotient identity]
\label{lem:compression}
For finite-dimensional $E$, the compression $T_E := P_E \Kc|_E$
has $\ip{T_E u}{u}/\|u\|^2 = \ip{\Kc u}{u}/\|u\|^2$ for
$u\in E\setminus\{0\}$.
\end{lemma}

%% ============================================================
\section{Gram Coercivity as Stability Input}
\label{sec:gram}
%% ============================================================

\begin{assumption}[Gram coercivity --- Paper~I]
\label{ass:gram}
$\lambda_{\min}(G_{c,N}) \geq \alpha_N(c) \asymp c^{-1/2}$
for all large $c$. (Proved unconditionally in Paper~I.)
\end{assumption}

\begin{proposition}[Gram coercivity does not imply spectral gaps]
\label{prop:gram-fails}
For any $\alpha > 0$ and any $N \geq 2$, there exists a self-adjoint
operator $T$ on an $N$-dimensional space with
$\lambda_{\min}(G) \geq \alpha$ (Gram matrix of a trial basis)
but all consecutive eigenvalue gaps equal to zero.
\end{proposition}

\begin{proof}
Let $\{e_j\}_{j=1}^N$ be any orthonormal basis and set
$T := \mu \cdot \mathrm{Id}$ for any $\mu > 0$.
Then the Gram matrix in basis $\{e_j\}$ is the identity,
so $\lambda_{\min}(G) = 1 \geq \alpha$ for $\alpha \leq 1$.
All $N$ eigenvalues of $T$ equal $\mu$, giving zero gaps.
For $\alpha > 1$, rescale: replace $e_j$ by $\alpha^{1/2} e_j$;
then $G = \alpha I$, $\lambda_{\min}(G) = \alpha$, and $T$ still
has all eigenvalues equal.
\end{proof}

\begin{remark}[What Gram coercivity does and does not give]
\label{rem:gram-limit}
Proposition~\ref{prop:gram-fails} is the key negative result.
Gram coercivity guarantees:
\begin{itemize}[nosep]
  \item \emph{Coordinate stability} (Gram invertibility);
  \item \emph{Rayleigh quotient reliability} (Lemma~\ref{lem:compression}).
\end{itemize}
It does \emph{not} guarantee any spectral gap, for any value of
$\lambda_{\min}(G)$.
Two additional inputs are required:
local spectral nondegeneracy (Assumption~\ref{ass:nondeg})
and spectral window isolation (Assumption~\ref{ass:window}).
\end{remark}

%% ============================================================
\section{Compressed Spectral Transfer}
\label{sec:transfer}
%% ============================================================

\subsection{Assumptions}

\begin{assumption}[Local spectral nondegeneracy]
\label{ass:nondeg}
For the compression $T_{c,N} := P_{E_{c,N}}\Kc|_{E_{c,N}}$,
there exists a basis of $E_{c,N}$ in which
\begin{equation}\label{eq:nondeg}
  \|T_{c,N} - \mathrm{diag}(\mu_j^{(c)})\|_{\op} \leq \delta(c),
\end{equation}
with $\min_j|\mu_j^{(c)} - \mu_{j+1}^{(c)}| \geq \delta_{\mathrm{model}}(c)$
and $\delta(c) = o(\delta_{\mathrm{model}}(c))$.
\end{assumption}

\begin{remark}[Operationalising Assumption~\ref{ass:nondeg}]
For $\Kc$-invariant spaces (e.g.\ $V_N$), \eqref{eq:nondeg} is trivially
satisfied with $\mu_j = \lam{c}{j}$ but then $\delta_{\mathrm{model}}$
is the gap itself (circular).
For a non-invariant trial space $\widetilde{V}_N$, verifying \eqref{eq:nondeg}
requires: (1) a microlocal normal form; (2) off-diagonal control via
noncommutativity of $\Pi_{\mathrm{space}}$ and $\Pi_{\mathrm{phase}}$;
(3) identification of the semiclassical scale of symbol separation.
This is the programme for Paper~XIV.
\end{remark}

\begin{assumption}[Spectral window isolation]
\label{ass:window}
The window $W_c := [\mu_{n+1}^{(c)} - \delta(c),\, \mu_n^{(c)} + \delta(c)]$
contains exactly two eigenvalues of $\Kc$,
namely $\lam{c}{n}$ and $\lam{c}{n+1}$.
\end{assumption}

\begin{remark}[Roles of the three assumptions]
\label{rem:three-roles}
\begin{center}
\renewcommand{\arraystretch}{1.2}
\begin{tabular}{lll}
\toprule
\textbf{Assumption} & \textbf{Type} & \textbf{Role} \\
\midrule
Gram coercivity (\ref{ass:gram}) & Stability & Coordinate non-collapse \\
Spectral nondegeneracy (\ref{ass:nondeg}) & Local spectral geometry & Forces compressed gap \\
Window isolation (\ref{ass:window}) & Index identification & Lifts to PSWF gap \\
\bottomrule
\end{tabular}
\end{center}
None of the three suffices alone.
Theorem~A is a \emph{lifting mechanism}: the depth lies entirely in
verifying the assumptions, not in the theorem itself.
\end{remark}

\subsection{The compressed-to-global gap lemma}

\begin{lemma}[Compressed-to-global gap transfer]
\label{lem:comptoglobal}
Let $A: H \to H$ be compact self-adjoint,
$E \subset H$ finite-dimensional with compression
$T_E = P_E A|_E$.
Let $\rho \geq 0$ and suppose:
\begin{enumerate}[label=(C\arabic*),nosep]
  \item $\mathrm{dist}(\sigma_j(T_E),\, \sigma(A)) \leq \rho$
    for $j = n, n+1$;\label{c1}
  \item $\sigma_n(T_E) - \sigma_{n+1}(T_E) \geq \Delta > 0$;\label{c2}
  \item The window $W := [\sigma_{n+1}(T_E) - \rho,\,
    \sigma_n(T_E) + \rho]$ contains exactly two eigenvalues of $A$,
    denoted $\lambda_p(A) \geq \lambda_{p+1}(A)$.\label{c3}
\end{enumerate}
Then
\begin{equation}\label{eq:comptoglobal}
  \lambda_p(A) - \lambda_{p+1}(A) \geq \Delta - 2\rho.
\end{equation}
If additionally $E$ contains vectors realising the min-max levels $p$
and $p+1$ for $A$ (i.e.\ $\lambda_p(A) = \lambda_p(T_E)$ holds with
respect to compatible min-max test spaces), then $p = n$.
\end{lemma}

\begin{proof}
By~\ref{c1}, each $\sigma_j(T_E)$ is within $\rho$ of some eigenvalue
of $A$; by~\ref{c3} the only eigenvalues of $A$ within $\rho$ of
$W$ are $\lambda_p$ and $\lambda_{p+1}$.
Hence $\lambda_p \in [\sigma_n(T_E) - \rho,\, \sigma_n(T_E) + \rho]$
and $\lambda_{p+1} \in [\sigma_{n+1}(T_E) - \rho,\, \sigma_{n+1}(T_E) + \rho]$.
Subtracting and using~\ref{c2}:
\[
  \lambda_p - \lambda_{p+1}
  \geq (\sigma_n - \sigma_{n+1}) - 2\rho \geq \Delta - 2\rho.
\]
\end{proof}

\begin{remark}[On condition~(C1) and the Weyl-perturbation bound]
\label{rem:C1}
In the setting of Theorem~A, (C1) follows directly from
Lemma~\ref{lem:compression} and the spectral mapping theorem:
compressed eigenvalues are Rayleigh quotients of $\Kc$, hence
$\mathrm{dist}(\sigma_j(T_{c,N}), \sigma(\Kc)) = 0$ exactly
(every Rayleigh quotient lies in the numerical range, and
for compact self-adjoint operators the numerical range closure
coincides with the convex hull of the spectrum).
Thus (C1) holds with $\rho = 0$ at this step;
the perturbation $\rho = \delta(c)$ enters only through the gap
between the compressed eigenvalues and the model values $\mu_j^{(c)}$
via Assumption~\ref{ass:nondeg}.
\end{remark}

\begin{remark}[On index identification]
\label{rem:index}
Lemma~\ref{lem:comptoglobal} establishes the gap between the two
eigenvalues of $A$ inside the window $W$, but does not in general
identify their global indices as $p = n$.
This index identification --- that the two window eigenvalues
are precisely $\lam{c}{n}$ and $\lam{c}{n+1}$ and not some
other consecutive pair --- is the content of
Assumption~\ref{ass:window} and is not derivable from the
compressed operator alone.
This is the correct locus of the difficulty: it is an external
spectral counting input, not a consequence of the transfer mechanism.
\end{remark}

\subsection{The theorem}

\begin{theorem}[Compressed Spectral Transfer --- Theorem~A]
\label{thm:A}
Under Assumptions~\ref{ass:gram}, \ref{ass:nondeg}, and~\ref{ass:window},
for all sufficiently large $c$ and relevant transition-window indices $n$:
\begin{equation}\label{eq:thmA}
  \gap{c}{n} \geq \tfrac{1}{2}\delta_{\mathrm{model}}(c).
\end{equation}
In particular, $\delta_{\mathrm{model}}(c) \asymp c^{-1/2}$
closes Gap-S and B-strong.
\end{theorem}

\begin{proof}
\textit{Step 1 (compressed gap via Weyl).}
By Assumption~\ref{ass:nondeg} and Weyl's inequality,
$|\sigma_j^{(c)} - \mu_j^{(c)}| \leq \delta(c)$, so
\[
  \sigma_n^{(c)} - \sigma_{n+1}^{(c)}
  \geq \delta_{\mathrm{model}}(c) - 2\delta(c) =: \Delta.
\]
Both $\sigma_n^{(c)}, \sigma_{n+1}^{(c)}$ lie in $W_c$.

\textit{Step 2 (compressed-to-global via Lemma~\ref{lem:comptoglobal}).}
Condition~(C1) holds with $\rho = 0$
(Remark~\ref{rem:C1}: compressed eigenvalues are exact Rayleigh quotients).
Condition~(C2) holds with $\Delta$ from Step~1.
Condition~(C3) is Assumption~\ref{ass:window}: $W_c$ contains exactly
$\lam{c}{n}$ and $\lam{c}{n+1}$.
Remark~\ref{rem:index} confirms that the index $p = n$ is supplied
by Assumption~\ref{ass:window}, not inferred from the compression.
Applying Lemma~\ref{lem:comptoglobal}:
$\gap{c}{n} \geq \Delta - 0 = \delta_{\mathrm{model}}(c) - 2\delta(c)$.

\textit{Step 3.}
$\delta(c) = o(\delta_{\mathrm{model}}(c))$ gives
$\gap{c}{n} \geq \tfrac{1}{2}\delta_{\mathrm{model}}(c)$.
\end{proof}

\begin{remark}[On the circularity for invariant trial spaces]
\label{rem:circular}
If $E_{c,N} = V_N$ ($\Kc$-invariant), $T_{c,N}$ is exactly diagonal
with eigenvalues $\lam{c}{0},\ldots,\lam{c}{N}$, so both
Assumption~\ref{ass:nondeg} (with $\mu_j = \lam{c}{j}$) and
Assumption~\ref{ass:window} reduce to Gap-S itself.
This is the precise mechanism by which Papers~X--XII fail:
every invariant-space approach inherits this circularity.
The only non-circular route requires a non-$\Kc$-invariant trial space
$\widetilde{V}_N$ for which the compressed symbol structure is
governed by the noncommutativity of $\Pi_{\mathrm{space}}$
and $\Pi_{\mathrm{frequency}}$.
\end{remark}

%% ============================================================
\section{Quasimode Transfer (the Solid Route)}
\label{sec:quasimode}
%% ============================================================

Theorem~B works directly with approximate eigenfunctions of $\Kc$
and is free of the circularity of Theorem~A.

\begin{definition}[Quasimode family]
$\{f_n^{(c)}\}$ is an $\varepsilon(c)$-quasimode family if
$\|(\Kc - \mu_n(c))f_n^{(c)}\| \leq \varepsilon(c)$, $\|f_n^{(c)}\| = 1$.
\end{definition}

\begin{assumption}[Local quasimode package]
\label{ass:quasimode}
\begin{enumerate}[label=(Q\arabic*),nosep]
  \item $\|f_n^{(c)}\| = 1$;
  \item $\|(\Kc - \mu_n(c))f_n^{(c)}\| \leq \varepsilon(c)$;
  \item $|\mu_n(c) - \mu_{n+1}(c)| \geq \delta_{\mathrm{model}}(c)$;
  \item $\|G_n^{(c)} - I\|_{\op} \leq \eta(c) \to 0$;
  \item \textbf{[Primary difficulty]}
    $[\mu_{n+1}(c) - \varepsilon(c),\, \mu_n(c) + \varepsilon(c)]$
    contains exactly two eigenvalues of $\Kc$.
\end{enumerate}
\end{assumption}

\begin{remark}[On the depth of (Q5)]
\label{rem:Q5}
(Q5) is a local spectral counting statement,
not a routine auxiliary condition.
It asserts $N(W) = 2$ for the transition window --- an upper and a lower
counting bound simultaneously.
This is essentially a local Widom-type law
in the transition regime $n \sim 2c/\pi$,
and is expected to encode most of the spectral rigidity of that regime.
It is the deepest open problem of this paper.
\end{remark}

\begin{lemma}[Kato--Temple gap transfer]
\label{lem:index}
Let $A$ compact self-adjoint, $\mu_n > \mu_{n+1}$,
$\varepsilon$-quasimodes $f_n, f_{n+1}$, $\|G - I\|_{\op} < 1$.
If $[\mu_{n+1} - \varepsilon,\, \mu_n + \varepsilon]$ contains exactly
two eigenvalues of $A$ (condition~(Q5)), these are $\lambda_n(A)$
and $\lambda_{n+1}(A)$ and
$\lambda_n(A) - \lambda_{n+1}(A) \geq (\mu_n - \mu_{n+1}) - 2\varepsilon$.
\end{lemma}

\begin{proof}
The spectral theorem gives $\mathrm{dist}(\mu_j, \sigma(A)) \leq \varepsilon$;
the Gram condition implies linear independence;
(Q5) provides the exact count and index identification;
ordering and subtracting gives the gap bound.
\end{proof}

\begin{theorem}[Quasimode Transfer --- Theorem~B]
\label{thm:B}
Under Assumption~\ref{ass:quasimode}:
$\gap{c}{n} \geq \delta_{\mathrm{model}}(c) - 2\varepsilon(c)$.
If $\varepsilon(c) = o(\delta_{\mathrm{model}}(c))$,
then $\gap{c}{n} \geq \frac{1}{2}\delta_{\mathrm{model}}(c)$.
\end{theorem}

\begin{proof}
Apply Lemma~\ref{lem:index} with $A = \Kc$.
\end{proof}

%% ============================================================
\section{Airy-Scale Conditional Consequence}
\label{sec:airy}
%% ============================================================

\begin{assumption}[Airy transition asymptotics]
\label{ass:airy}
Assumption~\ref{ass:quasimode} holds in the regime $n \sim 2c/\pi$
with $\delta_{\mathrm{model}}(c) \asymp c^{-1/3}$
and $\varepsilon(c) = o(c^{-1/3})$.
\end{assumption}

\begin{remark}[Status of Assumption~\ref{ass:airy}]
The scale $c^{-1/3}$ is expected from Airy quantization and
Widom--Deift transition asymptotics.
Proving it requires: (i) steepest-descent analysis of the PSWF
transition kernel; (ii) Airy quasimodes with residual $o(c^{-1/3})$;
(iii) condition~(Q5) in this regime.
None of these has been established from first principles.
\end{remark}

\begin{corollary}[Target theorem at Airy scale --- Corollary~C]
\label{cor:C}
Under Assumption~\ref{ass:airy},
$\gap{c}{n} \geq Cc^{-1/3}$ for all large $c$.
\end{corollary}

\begin{remark}[Corollary~C as a target theorem]
\label{rem:corC-status}
Corollary~C is a \emph{semiclassical prediction}, not a theorem.
It records the output one would obtain from a complete Airy quasimode
package and defines precisely what needs to be constructed.
\end{remark}

%% ============================================================
\section{Mechanism Table}
\label{sec:hierarchy}
%% ============================================================

\begin{center}
\renewcommand{\arraystretch}{1.35}
\begin{tabular}{lllp{3.2cm}}
\toprule
\textbf{Mechanism} & \textbf{Type} & \textbf{Output} & \textbf{Status} \\
\midrule
Gram coercivity (Paper~I)
  & Stability
  & Coordinate non-collapse
  & \checkmark\ unconditional \\
Prop.~\ref{prop:gram-fails}
  & Negative result
  & Gram coerc.\ $\not\Rightarrow$ gap
  & \checkmark\ proved \\
Comp.-to-global (Lem.~\ref{lem:comptoglobal})
  & Linear algebra
  & $\Delta_{\mathrm{comp}} - 2\rho \leq \gap{c}{n}$
  & \checkmark\ proved \\
Compressed nondeg.\ (Ass.~\ref{ass:nondeg})
  & Local spectral geometry
  & Compressed gap $\asymp c^{-1/2}$
  & open; non-circular only for $\widetilde{V}_N$ \\
Window isolation (Ass.~\ref{ass:window})
  & Index identification
  & $p = n$ in Lem.~\ref{lem:comptoglobal}
  & open (local counting) \\
Quasimode transfer (Thm~B)
  & Local spectral geometry
  & $\gap{c}{n} \geq \delta_{\mathrm{model}}$
  & conditional on Q1--Q5 \\
\textbf{(Q5)}
  & \textbf{Local spectral counting}
  & \textbf{Index identification for Thm~B}
  & \textbf{open; primary difficulty} \\
Airy quasimodes (Cor~C)
  & Semiclassical
  & $\gap{c}{n} \geq Cc^{-1/3}$
  & target theorem \\
\bottomrule
\end{tabular}
\end{center}

%% ============================================================
\section{Open Problems}
\label{sec:open}
%% ============================================================

\begin{problem}[Non-invariant trial space and compressed symbol analysis]
\label{prob:nondeg}
Construct $\widetilde{V}_N = \Pi_{\mathrm{phase}}\Pi_{\mathrm{space}} L^2$
and verify Assumption~\ref{ass:nondeg} with
$\delta_{\mathrm{model}} \asymp c^{-1/2}$:
(1) microlocal normal form;
(2) off-diagonal control via noncommutativity of
$\Pi_{\mathrm{space}}, \Pi_{\mathrm{phase}}$;
(3) semiclassical scale of symbol separation.
Programme for Paper~XIV.
\end{problem}

\begin{problem}[Spectral window isolation]
\label{prob:window}
Verify Assumption~\ref{ass:window} for the trial space of
OP~\ref{prob:nondeg}: local interlacing / counting for
the compressed spectrum.
\end{problem}

\begin{problem}[(Q5): Local spectral counting in the transition regime]
\label{prob:counting}
Prove condition~(Q5) in the regime $n \sim 2c/\pi$:
the spectral window contains exactly two eigenvalues of $\Kc$.
This is a local Widom-type counting problem and encodes most of
the transition-window spectral rigidity.
\emph{Deepest open problem in this paper.}
\end{problem}

\begin{problem}[Airy quasimode package]
\label{prob:quasimodes}
Construct quasimodes satisfying Assumption~\ref{ass:airy}
(residual $o(c^{-1/3})$, spacing $\asymp c^{-1/3}$,
condition~(Q5) in the Airy regime).
\end{problem}

\begin{problem}[Fold amplitude regime]
\label{prob:fold}
Determine $\alpha$ for the PSWF fold integral $\Phi_\alpha(u)$
(\texttt{fold\_model.tex}, \texttt{Waschtl904/prolate-primes-paper}).
\end{problem}

%% ============================================================
\section{Programme State After Paper~XIII}
\label{sec:state}
%% ============================================================

\begin{center}
\renewcommand{\arraystretch}{1.4}
\begin{tabular}{p{4.5cm} p{2.4cm} p{2.6cm} p{1.2cm}}
\toprule
\textbf{Claim} & \textbf{Source} & \textbf{Depends on} & \textbf{Status} \\
\midrule
Gram coerc.\ $\asymp c^{-1/2}$
  & Paper~I & --- & \checkmark \\
B-strong $\Leftrightarrow$ $\varepsilon_N \geq Cc^{-1/2}$
  & Paper~IX & --- & \checkmark \\
All direct routes fail
  & Papers~X--XII & --- & \checkmark \\
\textbf{Gram coerc.\ $\not\Rightarrow$ gap}
  & Prop.~\ref{prop:gram-fails} & --- & \textbf{\checkmark\ (new)} \\
\textbf{Compressed-to-global transfer}
  & Lem.~\ref{lem:comptoglobal} & --- & \textbf{\checkmark\ (new)} \\
\textbf{Circularity of inv.\ spaces}
  & Rem.~\ref{rem:circular} & --- & \textbf{\checkmark\ (new)} \\
\textbf{(Q5) = primary difficulty}
  & Rem.~\ref{rem:Q5} & --- & \textbf{\checkmark\ (new)} \\
Gap-S via Thm~A
  & This paper
  & Ass.~\ref{ass:gram}+\ref{ass:nondeg}+\ref{ass:window}
  & (cond.) \\
Gap-S via Thm~B
  & This paper & Q1--Q5 & (cond.) \\
Gap-S at $c^{-1/3}$ (Cor~C)
  & This paper & Airy pkg & (target) \\
Local spectral nondeg.
  & OP~\ref{prob:nondeg} & $\widetilde{V}_N$ & open \\
(Q5) local counting
  & OP~\ref{prob:counting} & Widom & open \\
Airy quasimodes
  & OP~\ref{prob:quasimodes} & --- & open \\
\bottomrule
\end{tabular}
\end{center}

\medskip\noindent
\textbf{Conceptual summary.}
\begin{quote}\itshape
Gap-S $\Longleftarrow$ local spectral rigidity / quasimode separation,
not Gap-S $\Longleftarrow$ Gram coercivity alone.
(Proposition~\ref{prop:gram-fails} proves the failure of the latter.)

\smallskip
Coordinate stability (Gram), local spectral geometry (nondegeneracy),
and index identification (window isolation / Q5) are logically independent.
All invariant-space approaches are circularly dependent on Gap-S
(Remark~\ref{rem:circular}).
The primary open problem is (Q5):
local spectral counting in the transition window.
\end{quote}

\begin{thebibliography}{99}
\bibitem{paper1} \textit{Paper~I: Gram Coercivity}, \texttt{paper1.tex}, 2026.
\bibitem{paper9} \textit{Paper~IX: B-strong}, \texttt{paper9\_superconcentration.tex}, 2026.
\bibitem{paper12} \textit{Paper~XII: Direct Coercivity}, \texttt{paper12\_direct\_coercivity.tex}, 2026.
\bibitem{foldmodel} \textit{Fold Amplitude}, \texttt{fold\_model.tex},
  \texttt{Waschtl904/prolate-primes-paper}, 2026.
\bibitem{Kato1966} T.~Kato, \textit{Perturbation Theory for Linear Operators},
  Springer, 1966.
\bibitem{Osipov2013} A.~Osipov, V.~Rokhlin, H.~Xiao,
  \textit{Prolate Spheroidal Wave Functions of Order Zero}, Springer, 2013.
\bibitem{Widom1964} H.~Widom, \textit{Asymptotic behavior of the eigenvalues
  of certain integral equations},
  Trans.\ Amer.\ Math.\ Soc.\ \textbf{109} (1964), 278--295.
\end{thebibliography}

\end{document}
